\documentclass[../../main.tex]{subfiles}
\graphicspath{{\subfix{../images/}}} % 指定图片目录，后续可以直接使用图片文件名。
\begin{document}
\section{2025年10月博士生资格考试原题}

\subsection{流体力学}
圆管中

\begin{enumerate}
    \item 列举涉及的重要物理量及其量纲
    \item 通过量纲分析流量 $Q\propto \Delta P$ [提示: $\begin{aligned}Q\propto \frac{\Delta P}{L}\end{aligned}$]
    \item 写出约化 navier-stokes 方程, 证明只需求解速度分布 $v(r)$
    \item 代入边界条件精确解 $v(r)$
    \item 精确解流量 $Q$
    \item 比较精确解和量纲分析结果的异同
\end{enumerate}

\subsection{应力学}

\begin{enumerate}
    \item 画图说明应力概念
    \item 解释为什么应力是张量
    \item 说明什么时候应力可被简化为压强
\end{enumerate}

\subsection{材料学}

\begin{enumerate}
    \item 解释概念 "杨氏模量"
    \item 排序花岗岩, 钢铁, 橡胶材料的杨氏模量大小
    \item 解释概念 "泊松比", 并且举例说明负泊松比材料
\end{enumerate}

\subsection{布朗运动}

\begin{enumerate}
    \item 解释概念 "布朗运动"
    \item 写出朗之万方程, 并对其进行变量解释
    \item 解释概念 "涨落耗散关系"
    \item 假定布朗粒子处于重力势中, 求解其在高度上的密度分布
\end{enumerate}

\subsection{聚合物}

\begin{enumerate}
    \item 解释概念 "理想链" 和 "非理想链", 并阐述两者之间的区别
    \item 求解三维空间中理想链的熵弹性力
    \item 解释概念 "Flory 理论"
    \item 若将聚合物溶解于某溶剂中, 阐述溶剂和温度对其平均半径的影响
\end{enumerate}

\subsection{液晶}

\begin{enumerate}
    \item 解释液晶的三种相: 丝状/向列相, 层状相和螺旋相
    \item 解释 splay, twist 和 bend 三种自由能
    \item 解释液晶的序参量定义
    \item 解释概念 "液晶的 $\begin{aligned}\pm\frac{1}{2}\end{aligned}$ 缺陷"
\end{enumerate}

\end{document}